\section{Background}

\subsection{Clinical and robotic motivation}

Technical skill matters clinically. In a study of bariatric surgeons whose operating videos were rated by peers, the lowest-rated quartile had markedly higher rates of complications, reoperation and readmission than the highest-rated quartile \citep{birkmeyer2013skill}. Skill is still assessed mainly by expert observers with structured checklists and global rating scales such as OSATS \citep{martin1997osats}. That process is slow, expensive and partly subjective, so it is hard to give trainees frequent feedback or to apply it consistently for credentialing \citep{reiley2011review}. Robot-assisted minimally invasive surgery changes the situation because the robot already records what the surgeon does. The da~Vinci system logs the pose, velocity and gripper angle of every manipulator at a fixed rate, and the endoscope records the operative field. Objective skill assessment can therefore be computed from data that are collected anyway, which makes it a natural problem for medical robotic intelligence. A model of this kind is only useful in practice if it works for a surgeon it has never seen. A trainee who is assessed by such a system will not have contributed trials to its training set, so the evaluation has to hold out whole surgeons rather than individual trials.

\subsection{Related work}

The JHU-ISI Gesture and Skill Assessment Working Set (JIGSAWS) is the standard public benchmark for this problem \citep{gao2014jigsaws}. It contains synchronised kinematics and stereo video of eight surgeons who each performed suturing, knot tying and needle passing on a bench-top model, together with a self-reported skill level and a modified OSATS score per trial. Benchmark protocols for the dataset distinguish leave-one-supertrial-out (LOSO), in which one repetition of every surgeon is held out, from leave-one-user-out (LOUO), in which all trials of one surgeon are held out \citep{ahmidi2017benchmarks}. Kinematic approaches fall roughly into three groups. The first summarises a trial with global motion statistics such as completion time, path length, economy of motion and the number of gripper movements, which are close to the metrics that expert raters use and are cheap to compute \citep{reiley2011review, lefor2020motion, anh2020features}. The second describes the texture of the motion signal. \citet{zia2018rmis}, the background reference of this topic, computed approximate entropy \citep{pincus1991apen} and related measures on every kinematic channel and reported 88.2\% accuracy for Suturing under LOUO with a nearest-neighbour classifier, and the same group applied the idea to video and accelerometer data \citep{zia2018video}. The third learns features end-to-end, for example with convolutional networks on the raw kinematic sequence \citep{fawaz2018cnn, wang2018cnn} or on video clips \citep{funke2019video, liu2021unified}, and recent video transformers decode surgeon activity and skill from clinical recordings \citep{kiyasseh2023vit, lavanchy2021threestage}. Many of the highest published numbers use LOSO, in which every test surgeon also appears in the training set, and accuracies near 100\% are common under that split \citep{fawaz2018cnn, funke2019video}. Under LOUO the task becomes generalisation to new people from only seven training surgeons, and the uncertainty of any estimate is large \citep{varoquaux2018cv}. Results obtained under the two splits are therefore not comparable. The class distribution adds a second difficulty. Novices contribute almost half of the Suturing trials, so accuracy rewards a classifier that favours the majority class, and the literature on imbalanced learning recommends class-balanced training together with metrics that weight every class equally, such as macro-F1 \citep{he2009imbalanced}. We follow both recommendations.

\subsection{Motivation for the approach}

The topic prescribes the core experiment: a kinematics-only baseline built from simple trial-level statistics, video features from sampled frames or short clips, a simple fusion of the two, and a comparison with the majority class under LOUO with macro-F1. Our design choices follow from two properties of the data. First, JIGSAWS Suturing has 39 usable trials from eight surgeons, and the skill label is constant within a surgeon. The effective sample size for learning what distinguishes skill levels is therefore closer to eight than to 39, with only two surgeons in each of the intermediate and expert classes. This favours low-capacity models with few free parameters and strong regularisation over end-to-end deep networks, and we use frozen pretrained image features with a linear classifier rather than any fine-tuning. Second, because labels belong to surgeons, any feature that identifies the surgeon also predicts the label inside the training set without capturing skill. We therefore treat shortcut detection as part of the method and add two measurements to the required experiment: a probe of how strongly each feature block encodes surgeon identity, and controls that measure how much of the decision is explained by the pace of the trial. As the focused improvement we test the central idea of the background reference, namely that expertise shows in the regularity of movement rather than only in its speed, by adding motion-texture features to the simple statistics. All models, metrics and comparisons were fixed in a written protocol before the first result was computed, so that the choice of the reported model could not be influenced by the outcome.
